\section{Introduction}
\label{sec:introduction}

Language-model agents can now read code, run experiments, and write prose with
little supervision. A large body of work has shown that structured
reasoning-and-acting loops~\cite{yao2023react}, self-reflection over past
attempts~\cite{shinn2023reflexion}, open-ended skill acquisition~\cite{wang2023voyager},
and purpose-built agent--computer interfaces for software
engineering~\cite{yang2024sweagent} materially raise what a single agent can do.
More recently, systems have attempted the full arc of scientific work ---
ideation, experimentation, and manuscript writing~\cite{lu2024aiscientist} --- and
external efforts have reported progress toward automating AI research
itself~\cite{recursive2026firststeps}. The frontier is no longer whether an agent
\emph{can} take a useful action; it is whether an agent can be trusted to take
thousands of them, unattended, without the run quietly falling apart.

That is the problem Argus targets. Give a capable model one hard objective and
enough wall-clock time, and the failure mode is rarely a failure of reasoning.
It is a failure of the \emph{execution substrate}. We have repeatedly observed
four distinct ways a long-horizon run degrades:

\begin{description}[leftmargin=1.4em, style=nextline]
  \item[High activity, low decision.] The agent stays busy --- re-reading the
    same files, rewriting the same status note, re-polling a healthy job ---
    while making no decision that moves the research forward. Token spend and
    round count climb; the actual frontier does not. Measuring ``progress'' by
    activity is exactly the trap.
  \item[Context amnesia.] A single mission's session is resumed dozens or
    hundreds of times. The underlying model compacts its context lossily on each
    overflow, discarding working memory, so the agent forgets what it already
    tried and re-treads dead ends. This is a structural property of unbounded
    session growth, not a prompt-quality problem.
  \item[Background-job babysitting.] The agent launches a long experiment --- a
    training run that already has its own health monitor --- and then blocks the
    main line of work polling it every round, instead of advancing independent
    tasks while it waits.
  \item[Premature acceptance.] The agent declares victory. Without an
    independent authority checking the claim against evidence, ``done'' becomes a
    self-report, and a self-report is not a result.
\end{description}

Layered on top of these is an \textbf{integrity} risk that is unique to
autonomous benchmarking: an agent optimizing a metric will, if permitted, exploit
the harness rather than solve the task --- hardcoding an answer, memorizing a
fixed evaluation input distribution, or editing a validator to report success.
An autonomous research system that cannot resist this is not merely unreliable;
it is actively misleading.

None of these five problems is solved by making the model smarter or the prompt
longer. They live in the substrate, so they must be fixed in the substrate. But
the substrate must fix them \emph{without} pretending to know better than the
agent about the research. A harness that starts guessing --- via keywords or
regexes --- whether an objective is ``really'' a paper task, whether an idea is
good enough, or whether a result should count, has simply moved the unreliability
somewhere harder to see. Argus's guiding constraint is therefore narrow and
strict: the harness may only enforce domain-agnostic invariants (budgets,
persistence, scheduling, structured I/O, and anti-cheat), and every scientific
judgment must be returned to the agent.

\begin{callout}[Scope of this report]
This is an engineering and design report at version 0.1. It documents what Argus
is built to do and how, grounded in the current source tree of the
\code{argus-skill} repository~\cite{argus_repo}. It reports evidence status
honestly and conservatively; where a program is ongoing, it says so, and it does
not convert an external baseline or a preliminary observation into an Argus
result.
\end{callout}

The remainder of the report proceeds as follows.
Section~\ref{sec:design} states the design philosophy that everything else
follows from. Section~\ref{sec:architecture} describes the four-role
architecture. Section~\ref{sec:reliability} details the reliability mechanisms.
Section~\ref{sec:evaluation} lays out the evidence tiers and measurement
discipline. Section~\ref{sec:programs} surveys the research programs and the
operator workbench. Section~\ref{sec:limitations} is candid about limitations and
responsible operation, and Section~\ref{sec:conclusion} closes with a roadmap.
